\usection{Conclusiones}

El presente trabajo de grado tuvo como propósito el desarrollo de NodeLand, un lenguaje de programación visual basado en flujos de datos orientado a fomentar el pensamiento computacional en niños y jóvenes de 8 a 16 años. A partir de los resultados obtenidos en cada etapa del proceso, se derivan las siguientes conclusiones.

En relación con el análisis de la problemática y la caracterización del lenguaje, se concluye que el paradigma de flujo de datos (dataflow) constituye una alternativa teórica y pedagógicamente viable en comparación con el flujo de control (control flow) para la introducción de conceptos computacionales avanzados en etapas tempranas de formación. El análisis demostró que las herramientas educativas existentes, al estar fundamentadas exclusivamente en la ejecución secuencial, tienden a formar modelos mentales rígidos que dificultan la posterior comprensión de la concurrencia y los sistemas distribuidos. Asimismo, la revisión de las etapas de desarrollo cognitivo de Piaget permitió constatar que la representación visual explícita de las dependencias de datos, propia del paradigma dataflow, resulta compatible con el pensamiento concreto-operacional, en contraste con la abstracción implícita del contador de programa imperativo, lo que valida la pertinencia del paradigma seleccionado para el público objetivo.

En cuanto al diseño de la sintaxis visual y la semántica del lenguaje, se evidenció que la aplicación rigurosa de la Teoría de la Física de las Notaciones demostró ser una condición importante para el diseño de notaciones visuales dirigidas a niños. La codificación dual, la discriminabilidad perceptual por color y forma, y la correspondencia entre la representación visual y la semántica de ejecución redujeron efectivamente la carga cognitiva durante las pruebas con el grupo objetivo. Por otra parte, el diseño de una máquina de estados finitos para modelar el ciclo de vida de cada nodo reveló que la retroalimentación visual continua no es un complemento estético del entorno, sino un mecanismo fundamental de depuración pedagógica que permite al estudiante razonar sobre el flujo de la información sin necesidad de conocimiento abstracto previo.

Adicionalmente, el diseño arquitectónico del motor de ejecución permite derivar una conclusión de especial relevancia técnica y pedagógica. La decisión de fundamentar el motor sobre el principio de activación por disponibilidad de datos (data readiness), en lugar de un contador de programa secuencial, demostró que la concurrencia implícita no requiere ser gestionada explícitamente por el programador para ser real y observable. Múltiples nodos pueden activarse simultáneamente como consecuencia natural de la disponibilidad de sus entradas, sin que el estudiante deba invocar ningún mecanismo de sincronización o manejo de hilos. Esto implica que la concurrencia, lejos de ser un concepto avanzado que debe introducirse tardíamente, puede ser la condición por defecto del entorno de aprendizaje desde el primer programa construido por el usuario.

Respecto a la implementación del lenguaje y su entorno, los resultados permiten establecer que las tecnologías web estándar actuales son suficientes para materializar un entorno de programación distribuido, concurrente y ejecutable nativamente en el navegador, sin comprometer la experiencia del usuario ni la expresividad del lenguaje. El desarrollo incremental evidenció, además, que la ausencia de soporte para recursividad en las primeras versiones del motor no fue un error de diseño, sino una limitación inherente a los lenguajes dataflow puros. Su resolución, mediante la instanciación controlada de subflujos, demostró que es posible alcanzar la completitud de Turing dentro del paradigma sin recurrir a construcciones imperativas.

En lo que concierne a la validación del lenguaje desarrollado, se verificó que NodeLand satisface simultáneamente la pertinencia en contextos educativos, y la usabilidad adecuada para el público objetivo. La evaluación con expertos confirmó que el enfoque visual basado en dependencias de datos facilita la tangibilización de conceptos matemáticos y lógicos abstractos, ampliando el alcance potencial de la herramienta más allá de la computación hacia otras áreas del currículo escolar. Por su parte, las pruebas con niños y jóvenes demostraron que la representación de un programa como un grafo de conexiones lógicas es asimilable de manera intuitiva, lo que contradice la suposición implícita de que los paradigmas distintos al flujo de control son inherentemente más difíciles para usuarios novatos, evidenciando incluso indicios prometedores de que la dificultad históricamente atribuida a la programación concurrente radica, al menos en parte, en la notación empleada y no necesariamente en la naturaleza del concepto.

Finalmente, respecto a la elaboración de la documentación formal del lenguaje, se puede afirmar que el manual de usuario desarrollado cumple una función que excede la de un recurso de referencia técnica. Al estar orientado al docente, sirve de puente entre las capacidades del lenguaje y su aprovechamiento en el aula. La documentación realizada habilita al educador para diseñar ciclos de resolución de problemas progresivos y proponer desafíos de mayor complejidad que los contemplados en el módulo de desafíos integrado, adaptando NodeLand a los contenidos y objetivos curriculares de su contexto particular. Esto implica que la calidad y completitud del manual determina, en gran medida, el alcance real del impacto pedagógico de la herramienta, dado que un lenguaje cuyas capacidades no son correctamente comprendidas por el docente difícilmente será aprovechado más allá de sus usos más superficiales.

Complementariamente, la construcción de un ecosistema de recursos del lenguaje, materializado en la plataforma social de proyectos compartidos y el sistema de extensiones, ofrece al docente un repertorio de ejemplos concretos sobre los cuales fundamentar sus propuestas didácticas, enriqueciendo la documentación del sistema con evidencia práctica generada por la propia comunidad de usuarios, en un ciclo de retroalimentación constante que favorece a todos los involucrados.